% !TeX root = main.tex
\section{Introduction}
\label{s:introduction}

Exploiting visual sensors that capture depth information, research has recently started exploring the capture, transmission, and display of \textit{volumetric} videos.
%
Each frame of a volumetric video captures a scene or an object in three dimensions.
%
A frame is often represented as a point cloud, a collection of points on surfaces of a person or object, together with each point's position (also known as \textit{geometry}) and color.

More than regular (2D) videos, volumetric videos stress compute and communication, since point clouds can be large.
%
They differ from regular (2D) videos in one other aspect.
%
Users can \textit{interact} with them with 6 degrees of freedom, by viewing a frame from any angle, at any location in the frame.
%
Users can shift perspective by moving in a designated space while wearing a headset that tracks their movement, and renders each frame using the current perspective.
%; the headset can then render the current frame using the current perspective.

Volumetric video promises a degree of immersion comparable to physical presence.
%
Imagine live two-way streaming of entire physical spaces
(\textit{e.g.}, areas with multiple participants and objects in the environment) between two locations.
%
Users can virtually navigate each other's physical spaces live, with these movements being instantly visible at the remote end.
%
However, complete physical scene descriptions are large (upwards of 10~MB per video frame, \cref{s:background}).
%
Fortunately, average broadband bandwidths worldwide are now about 100~Mbps, so full-scene volumetric video transmission is within the realm of feasibility.
%
Even with this, however, it will not be possible to transmit volumetric video without compression.

% Early research in this area explored on-demand streaming of volumetric video~\cite{groot,vivo,others}.
% %
% More recent work has considered \emph{live} streaming volumetric videos to enable, for example, two-way conferencing~\cite{starline,metastream}.
% %
% However, this work has not yet exploited the full power of 6DoF interactivity; instead, most of it has focused on transmitting volumetric videos of a single participant, or of their face or torso.
% %
% If, instead, it was possible to live stream entire physical spaces
% (\textit{e.g.}, rooms, small performance areas with multiple participants and objects in the environment), users can experience a degree of immersion comparable to physical presence.

Point cloud compression is in its infancy, relative to 2D video compression.
%
State-of-the-art point cloud compression techniques are slow, often requiring hundreds of milliseconds to a few seconds to compress a single point cloud frame.
%
This renders them unusable for two-way live streaming applications like conferencing, which require low end-to-end latency.
%
Furthermore, relative to 2D video compression, inter-frame compression is less well developed, resulting in lower compression efficiency.
%
Finally, most point cloud compression codecs are not \textit{rate-adaptive}, \ie they cannot compress a point cloud to a given target bandwidth.
%
As such, volumetric video transmissions cannot be \textit{bandwidth-adaptive}~---~use available bandwidth to deliver high-quality video when possible, and to minimize \textit{stalls} when there are dips in available bandwidth due to contention.

\parab{Contributions.}
%
In this paper, we describe the design and implementation of \sysname, which transmits full-scene volumetric video at full frame rate (\ie 30 frames per second), achieves low end-to-end latency comparable to today's 2D video conferencing systems, and adapts to available bandwidth.
%
Rather than relying on point cloud compression, \sysname's key insight is that it is possible (and desirable) to re-use the rather extensive infrastructure developed for 2D video conferencing: efficient compression algorithms, rate-adaptive codecs that can encode and decode at full frame rate, and protocols like WebRTC for transmission of video (\cref{s:design}).
%
To this end, it develops novel, cheap techniques to encode the output of an array of RGB-D cameras, which can produce 3D scene representations, into \textit{two} 2-D video streams, one for geometry and the other for color.
%
This allows it to re-use 2D video technology, but by itself is not sufficient.
%
\sysname also develops a lightweight technique to track the receiver's perspective and \textit{cull} a point cloud to remove points not within the receiver's field of view.
%
Because humans are more sensitive to errors in geometry than color, \sysname
carefully apportions available bandwidth between the two video streams, resulting in high-quality, bandwidth-adaptive video delivery.

%\rn{I wonder if there is a way to sell this a bit more as challenges we had to overcome beyond the key insight? Also, is it obvious that volumetric video point clouds are generated by fusing RGBD feeds?}

Using a complete implementation of \sysname, which we plan to release publicly, we have conducted a user study and measured objective 3D quality metrics.
%
\sysname has significantly better visual quality than a bandwidth-adaptive \textit{oracle} of Draco~\cite{draco}, a popular point cloud compression technique.
%
It also has about 30\% higher quality than an approach that mimics a bandwidth-adaptive version of Project Starline~\cite{starline}.
%
These differences arise primarily due to the two key design choices: using efficient, bandwidth-adaptive 2D video compression, and culling views based on receiver perspective.
%
Experiments also demonstrate that \sysname can achieve about 300~ms end-to-end latency at 30 frames per second, comparable to the performance of existing 2D conferencing platforms. %We will release \sysname{} and our datasets post publication.

%\sanjay{Move para below earlier, maybe above Contributions, and move the "to our knowledge line as first line in Contributions?"}
%\ramesh{Sanjay, hope it is OK to keep it here, since introducing on-demand earlier will break the flow.}
\parab{Related work.}
%
To our knowledge, relative to prior work, \sysname is the first to explore full-scene, bandwidth-adaptive, low latency live streaming at full frame rate (\cref{tab:comparison}, \cref{s:related}).
%
Early research in this area has explored on-demand streaming of volumetric video~\cite{groot,vivo}.
%
Much of this uses point-cloud compression and optimizes it for speed.
%
More recent work has considered \emph{live} streaming volumetric videos to enable, for example, one-way~\cite{metastream} and two-way conferencing~\cite{starline}, but transmits constrained views (of a single participant, or of their face or torso).
%
None of this work has considered bandwidth-adaptive, full-frame transmission, though some of it~\cite{starline} uses  2D videos to represent the point cloud but compresses them at fixed quality.

%\rn{can we add `meeting the targets of existing 2D conferencing platforms'?}

